%=============================================================================
% COMPREHENSIVE TECHNICAL DOCUMENTATION
% Integrated Stereo SLAM System with Reinforcement Learning
% A-to-Z Line-by-Line Explanation with Deep Core Mathematics
%=============================================================================

\documentclass[12pt,a4paper,twoside]{report}

%----------------------------- PACKAGES -------------------------------------
\usepackage[utf8]{inputenc}
\usepackage[T1]{fontenc}
\usepackage[english]{babel}
\usepackage{geometry}
\usepackage{amsmath,amssymb,amsfonts,amsthm}
\usepackage{mathtools}
\usepackage{bm}               % bold math
\usepackage{graphicx}
\usepackage{xcolor}
\usepackage{listings}
\usepackage{fancyhdr}
\usepackage{hyperref}
\usepackage{tcolorbox}
\usepackage{booktabs}
\usepackage{array}
\usepackage{longtable}
\usepackage{multirow}
\usepackage{tikz}
\usepackage{algorithm}
\usepackage{algpseudocode}
\usepackage{enumitem}
\usepackage{caption}
\usepackage{subcaption}
\usepackage{float}
\usepackage{titlesec}
\usepackage{microtype}
\usepackage{lmodern}
\usepackage{setspace}
\usepackage{mdframed}
\usepackage{soul}
\usepackage{upquote}

%----------------------------- GEOMETRY ------------------------------------
\geometry{
  left=2.5cm, right=2.5cm,
  top=3cm, bottom=3cm,
  headheight=15pt
}

%----------------------------- COLORS --------------------------------------
\definecolor{codegreen}{rgb}{0,0.6,0}
\definecolor{codegray}{rgb}{0.5,0.5,0.5}
\definecolor{codepurple}{rgb}{0.58,0,0.82}
\definecolor{codeblue}{rgb}{0.13,0.29,0.53}
\definecolor{backcolour}{rgb}{0.97,0.97,0.97}
\definecolor{mathblue}{rgb}{0.05,0.2,0.55}
\definecolor{notegreen}{rgb}{0.0,0.4,0.0}
\definecolor{warnorange}{rgb}{0.7,0.35,0.0}
\definecolor{titleblue}{rgb}{0.1,0.1,0.55}

%----------------------------- LISTINGS (CODE) ----------------------------
\lstdefinestyle{pythonstyle}{
  backgroundcolor=\color{backcolour},
  commentstyle=\color{codegreen}\itshape,
  keywordstyle=\color{codeblue}\bfseries,
  numberstyle=\tiny\color{codegray},
  stringstyle=\color{codepurple},
  basicstyle=\ttfamily\footnotesize,
  breakatwhitespace=false,
  breaklines=true,
  captionpos=b,
  keepspaces=true,
  numbers=left,
  numbersep=5pt,
  showspaces=false,
  showstringspaces=false,
  showtabs=false,
  tabsize=4,
  frame=single,
  rulecolor=\color{codegray},
  xleftmargin=1em,
  xrightmargin=0.5em
}
\lstset{style=pythonstyle, language=Python}

%----------------------------- TCOLORBOX STYLES ---------------------------
\tcbuselibrary{skins,breakable,listings}

\newtcolorbox{mathbox}[1][]{
  colback=blue!5!white, colframe=blue!40!black,
  fonttitle=\bfseries, title={#1},
  breakable
}

\newtcolorbox{notebox}[1][Note]{
  colback=green!5!white, colframe=green!40!black,
  fonttitle=\bfseries, title={#1},
  breakable
}

\newtcolorbox{warnbox}[1][Warning]{
  colback=orange!5!white, colframe=orange!60!black,
  fonttitle=\bfseries, title={#1},
  breakable
}

\newtcolorbox{defbox}[1][Definition]{
  colback=gray!5!white, colframe=gray!50!black,
  fonttitle=\bfseries, title={#1},
  breakable
}

%----------------------------- HEADERS/FOOTERS ---------------------------
\pagestyle{fancy}
\fancyhf{}
\fancyhead[LE,RO]{\nouppercase{\leftmark}}
\fancyhead[RE,LO]{\textit{Integrated SLAM + RL Documentation}}
\fancyfoot[C]{\thepage}
\renewcommand{\headrulewidth}{0.4pt}

%----------------------------- HYPERREF ----------------------------------
\hypersetup{
  colorlinks=true,
  linkcolor=titleblue,
  filecolor=magenta,
  urlcolor=cyan,
  pdftitle={Integrated Stereo SLAM + RL: Complete Technical Manual},
  pdfauthor={Auto-Generated Documentation},
  pdfsubject={SLAM, Reinforcement Learning, Computer Vision},
}

%----------------------------- THEOREM ENVIRONMENTS ----------------------
\theoremstyle{definition}
\newtheorem{definition}{Definition}[chapter]
\newtheorem{theorem}{Theorem}[chapter]
\newtheorem{lemma}[theorem]{Lemma}
\newtheorem{proposition}[theorem]{Proposition}
\newtheorem{corollary}[theorem]{Corollary}
\newtheorem{remark}{Remark}[chapter]
\newtheorem{example}{Example}[chapter]

%----------------------------- CUSTOM COMMANDS ---------------------------
\newcommand{\vect}[1]{\bm{#1}}
\newcommand{\mat}[1]{\mathbf{#1}}
\newcommand{\R}{\mathbb{R}}
\newcommand{\E}{\mathbb{E}}
\newcommand{\prob}{\mathbb{P}}
\newcommand{\argmax}{\operatorname*{arg\,max}}
\newcommand{\argmin}{\operatorname*{arg\,min}}
\newcommand{\norm}[1]{\left\lVert #1 \right\rVert}
\newcommand{\abs}[1]{\left| #1 \right|}
\newcommand{\T}{\top}           % transpose
\newcommand{\SE}{\mathrm{SE}(3)}
\newcommand{\SO}{\mathrm{SO}(3)}
\newcommand{\code}[1]{\texttt{#1}}
\newcommand{\file}[1]{\texttt{\textcolor{codeblue}{#1}}}
\newcommand{\lineno}[1]{\textcolor{codegray}{\small[L#1]}}

%=============================================================================
%                            DOCUMENT BEGINS
%=============================================================================
\begin{document}

%----------------------------- TITLE PAGE --------------------------------
\begin{titlepage}
  \centering
  \vspace*{1.5cm}
  {\Huge\bfseries\color{titleblue}
    Integrated Stereo SLAM System\\[0.4em]
    with Reinforcement Learning\\[0.4em]
    Keyframe Selection\par}
  \vspace{1.5cm}
  {\Large\itshape A Complete A-to-Z Technical Manual\\
   Line-by-Line Code Analysis with Deep Mathematical Foundations\par}
  \vspace{2cm}
  \rule{\linewidth}{1pt}
  \vspace{0.5cm}
  {\large
    \begin{tabular}{ll}
      \textbf{Project:}   & Integrated Stereo SLAM + RL \\
      \textbf{Language:}  & Python 3.x \\
      \textbf{Modules:}   & OpenCV, NumPy, PyTorch, psutil \\
      \textbf{Domain:}    & Computer Vision, Robotics, RL \\
      \textbf{Doc Type:}  & Full Technical Reference Manual \\
    \end{tabular}
  }
  \vspace{0.5cm}
  \rule{\linewidth}{1pt}
  \vfill
  {\large\today}
\end{titlepage}

%----------------------------- TABLE OF CONTENTS -------------------------
\tableofcontents
\newpage
\listoffigures
\newpage
\listoftables
\newpage

%=============================================================================
\chapter*{Preface}
\addcontentsline{toc}{chapter}{Preface}
%=============================================================================

This document is a complete, exhaustive technical reference for the
\textbf{Integrated Stereo SLAM System with Reinforcement Learning-based Keyframe
Selection}. It covers every aspect of the project:

\begin{itemize}[leftmargin=1.5em]
  \item \textbf{Architecture overview}: how all modules fit together.
  \item \textbf{Line-by-line code explanation}: every non-trivial line of Python
        source code is annotated and explained.
  \item \textbf{Deep mathematics}: the linear algebra, probability theory,
        optimization, and RL theory behind every algorithm.
  \item \textbf{Data-flow diagrams}: how data moves from camera pixels to the
        3-D map and RL reward.
  \item \textbf{Configuration tables}: hyper-parameters and their effect.
\end{itemize}

\noindent The reader is assumed to have undergraduate-level knowledge of linear
algebra and basic Python. No prior robotics or machine-learning experience is
required; all concepts are introduced from first principles.

\setcounter{page}{1}

%=============================================================================
\chapter{Project Overview and Architecture}
%=============================================================================

\section{High-Level Purpose}

The project implements a real-time \textbf{Stereo Simultaneous Localization and
Mapping (SLAM)} pipeline for a mobile robot or wearable device equipped with two
IP cameras. It addresses a fundamental tension in battery-powered robotics:

\begin{center}
\emph{``High-accuracy localization requires many keyframes (expensive battery-wise);\\
      long battery life requires skipping frames (accuracy-costly).''}
\end{center}

\noindent A \textbf{Reinforcement Learning (RL)} agent observes the current
localization uncertainty and the battery state, then decides -- for every
incoming stereo frame -- whether to commit it as a keyframe or discard it.

\section{Repository File Map}

\begin{longtable}{|p{5.5cm}|p{8cm}|}
\caption{All project files and their roles}\label{tab:filemap}\\
\hline
\textbf{File path} & \textbf{Role} \\ \hline
\endfirsthead
\hline
\textbf{File path} & \textbf{Role} \\ \hline
\endhead
\file{integrated\_slam.py} (root) & Entry point: battery-aware SLAM without RL \\ \hline
\file{app/backend/camera\_stream\_handler.py} & MJPEG stream decoder; stereo camera manager \\ \hline
\file{app/backend/stereo\_vision.py} & Depth estimation, 3-D triangulation, Kabsch pose \\ \hline
\file{app/backend/advanced\_features.py} & SIFT/ORB/AKAZE feature extraction; loop closure; cloud export \\ \hline
\file{app/backend/visual\_odometry.py} & Frame-to-frame motion estimation; trajectory accumulation \\ \hline
\file{app/backend/stereo\_slam.py} & Full SLAM system: keyframing, loop closure, bundle adjustment \\ \hline
\file{app/backend/rl\_keyframe\_selector.py} & Q-learning agent; uncertainty/energy models \\ \hline
\file{app/backend/integrated\_slam\_rl.py} & Ultimate SLAM+RL entry point with full UI \\ \hline
\file{app/backend/dqn\_agent.py} & Deep Q-Network (PyTorch) for advanced keyframe selection \\ \hline
\file{app/backend/train\_rl\_agent.py} & RL training script with SLAM simulator \\ \hline
\file{app/backend/energy\_battery\_model.py} & Energy cost table; battery monitor; reward calculator \\ \hline
\file{app/backend/degradation\_policy.py} & 4-level graceful degradation of SLAM quality \\ \hline
\file{tools/launch\_*.py} & Launcher scripts for each mode \\ \hline
\file{data/*.json} & Saved SLAM maps (trajectory + point cloud) \\ \hline
\end{longtable}

\section{System Architecture Diagram}

\begin{figure}[H]
\centering
\begin{tikzpicture}[
  node distance=1.0cm and 1.5cm,
  every node/.style={draw, rounded corners, align=center, font=\small},
  arr/.style={->, thick, >=stealth}
]
% Row 1: Cameras
\node[fill=blue!15] (lc) {Left Camera\\(MJPEG URL)};
\node[fill=blue!15, right=of lc] (rc) {Right Camera\\(MJPEG URL)};
% Row 2: Stream handler
\node[fill=cyan!15, below=of lc, xshift=1.8cm] (sh) {StereoCameraManager\\camera\_stream\_handler.py};
% Row 3: Feature extraction
\node[fill=yellow!20, below=of sh] (fe) {AdvancedFeatureExtractor\\advanced\_features.py};
% Row 4: Branches
\node[fill=orange!20, below left=of fe] (vo) {VisualOdometryEngine\\visual\_odometry.py};
\node[fill=orange!20, below right=of fe] (sv) {StereoCalibratedPipeline\\stereo\_vision.py};
% Row 5: SLAM
\node[fill=red!15, below=of vo, xshift=1.8cm] (slam) {StereoSLAMSystem\\stereo\_slam.py};
% Row 6: RL
\node[fill=green!15, right=4cm of slam] (rl) {RLKeyframeSelector\\rl\_keyframe\_selector.py};
% Row 7: Battery
\node[fill=purple!15, above=of rl] (batt) {BatteryManager\\energy\_battery\_model.py};
% Row 8: Output
\node[fill=gray!15, below=1.5cm of slam] (out) {3-D Map + Trajectory\\(JSON / PLY)};

\draw[arr] (lc) -- (sh);
\draw[arr] (rc) -- (sh);
\draw[arr] (sh) -- (fe);
\draw[arr] (fe) -- (vo);
\draw[arr] (fe) -- (sv);
\draw[arr] (vo) -- (slam);
\draw[arr] (sv) -- (slam);
\draw[arr] (rl) -- node[draw=none,fill=none,sloped,above]{\tiny decide} (slam);
\draw[arr] (batt) -- (rl);
\draw[arr] (slam) -- (out);
\draw[arr] (slam) |- node[draw=none,fill=none,near start,right]{\tiny uncertainty} (rl);
\end{tikzpicture}
\caption{High-level data-flow of the integrated SLAM+RL system.}
\label{fig:arch}
\end{figure}

\section{Mathematical Notation Reference}

\begin{table}[H]
\centering
\caption{Symbols used throughout this document}\label{tab:notation}
\begin{tabular}{clp{7cm}}
\toprule
\textbf{Symbol} & \textbf{Type} & \textbf{Meaning} \\ \midrule
$\mat{T} \in \SE$ & $4\times4$ matrix & Rigid-body transformation (pose) \\
$\mat{R} \in \SO$ & $3\times3$ matrix & Rotation matrix \\
$\vect{t} \in \R^3$ & vector & Translation vector \\
$\mat{K}$           & $3\times3$ matrix & Camera intrinsic matrix \\
$f$                 & scalar & Focal length (pixels) \\
$b$                 & scalar & Stereo baseline (metres) \\
$d$                 & scalar & Disparity (pixels) \\
$Z$                 & scalar & Depth (metres) \\
$\mat{F}$           & $3\times3$ matrix & Fundamental matrix \\
$\mat{E}$           & $3\times3$ matrix & Essential matrix \\
$Q(s,a)$            & scalar & Q-value for state $s$, action $a$ \\
$\pi(s)$            & action & Policy mapping state to action \\
$\varepsilon$       & scalar & Exploration rate (epsilon-greedy) \\
$\gamma$            & scalar $\in[0,1]$ & Discount factor \\
$\alpha$            & scalar & Learning rate \\
$r$                 & scalar & Reward signal \\
\bottomrule
\end{tabular}
\end{table}

%=============================================================================
\chapter{Mathematical Prerequisites}
%=============================================================================

\section{3-D Rigid-Body Transformations}

\begin{definition}[Special Euclidean Group $\SE$]
$\SE$ is the group of rigid-body transformations in three dimensions:
\[
  \SE = \left\{ \mat{T} = \begin{bmatrix} \mat{R} & \vect{t} \\ \vect{0}^\T & 1 \end{bmatrix}
         \;\Bigg|\; \mat{R} \in \SO,\; \vect{t} \in \R^3 \right\}
\]
where $\SO = \{\mat{R}\in\R^{3\times3} \mid \mat{R}^\T\mat{R}=\mat{I},\;\det\mat{R}=+1\}$.
\end{definition}

\noindent A point $\vect{p}_w \in \R^3$ in world coordinates is transformed to
camera coordinates by:
\[
  \vect{p}_c = \mat{R}\,\vect{p}_w + \vect{t}
\]
In homogeneous form:
\[
  \tilde{\vect{p}}_c = \mat{T}\,\tilde{\vect{p}}_w, \quad
  \tilde{\vect{p}} = \begin{bmatrix} \vect{p} \\ 1 \end{bmatrix}
\]

\subsection{Pose Composition}

Given two consecutive poses $\mat{T}_1$ (world-to-camera-1) and
$\mat{T}_{12}$ (camera-1-to-camera-2), the cumulative pose is:
\[
  \mat{T}_2 = \mat{T}_1 \cdot \mat{T}_{12}
\]
This is exploited in \file{visual\_odometry.py} line 265:
\begin{lstlisting}
self.camera_pose = self.camera_pose @ T
\end{lstlisting}

\section{Camera Projection}

\begin{definition}[Perspective Projection]
A 3-D point $[X,Y,Z]^\T$ (camera frame) projects to image pixel $[u,v]^\T$:
\[
  u = \frac{f_x\,X}{Z} + c_x, \qquad v = \frac{f_y\,Y}{Z} + c_y
\]
where $(f_x,f_y)$ are focal lengths and $(c_x,c_y)$ is the principal point.
In matrix form with intrinsic matrix $\mat{K}$:
\[
  \mat{K} = \begin{bmatrix} f_x & 0 & c_x \\ 0 & f_y & c_y \\ 0 & 0 & 1 \end{bmatrix}
\]
\[
  \lambda\begin{bmatrix}u\\v\\1\end{bmatrix} = \mat{K}\begin{bmatrix}X\\Y\\Z\end{bmatrix}
\]
\end{definition}

\noindent The project approximates $f_x=f_y=f=500$ px, $c_x=320$, $c_y=240$
(file \file{stereo\_vision.py}, lines 41--42).

\section{Stereo Depth Estimation}

For a rectified stereo camera pair with baseline $b$ and focal length $f$,
the disparity $d = u_L - u_R$ relates to depth $Z$ as:
\[
  Z = \frac{f \cdot b}{d}
\]

\begin{mathbox}[Stereo Depth Formula]
\[
  Z = \frac{f \cdot b}{d}, \qquad X = \frac{(u-c_x)\cdot Z}{f}, \qquad Y = \frac{(v-c_y)\cdot Z}{f}
\]
These three equations triangulate a 3-D point from a single stereo match.
\end{mathbox}

\noindent In code (\file{stereo\_vision.py} lines 115--168):
\begin{lstlisting}
# OpenCV block matching disparity
disparity = self.stereo_matcher.compute(left_gray, right_gray)
               .astype(np.float32)

# Division by 16 converts StereoBM fixed-point output to real disparity
depth = (self.focal_length * self.baseline) / (disparity / 16.0)

# Back-projection to 3-D
X = (x_l - cx) * depth / self.focal_length   # horizontal
Y = (y_l - cy) * depth / self.focal_length   # vertical
Z = depth                                      # optical axis
\end{lstlisting}

\noindent \textbf{Why divide by 16?} OpenCV's \code{StereoBM} stores disparity as
a 16-bit fixed-point value with 4 fractional bits, so the raw integer value is
$16\times d_{\text{real}}$.  Dividing by 16 recovers the true sub-pixel disparity.

\section{Essential Matrix and Pose Recovery}

\begin{definition}[Essential Matrix]
Let $\vect{x}_1, \vect{x}_2$ be normalized image coordinates of the same scene
point in two views with relative rotation $\mat{R}$ and translation $\vect{t}$.
The \emph{essential matrix} $\mat{E}$ satisfies the epipolar constraint:
\[
  \vect{x}_2^\T \mat{E} \vect{x}_1 = 0, \qquad \mat{E} = [\vect{t}]_\times \mat{R}
\]
where $[\vect{t}]_\times$ is the skew-symmetric matrix of $\vect{t}$.
\end{definition}

\noindent Steps in \file{visual\_odometry.py} (lines 100--130):
\begin{enumerate}
  \item \textbf{Compute $\mat{E}$} using OpenCV RANSAC:
        \code{cv2.findEssentialMat(..., method=cv2.RANSAC, prob=0.999, threshold=1.0)}
  \item \textbf{Recover pose} $(R,t)$ via SVD decomposition of $\mat{E}$:
        \code{cv2.recoverPose(E, pts1, pts2)}
  \item \textbf{Update} cumulative pose: $\mat{T} \leftarrow \mat{T} \cdot \mat{T}_{new}$
\end{enumerate}

\begin{mathbox}[Essential Matrix SVD Decomposition]
Given $\mat{E} = \mat{U}\mat{\Sigma}\mat{V}^\T$ (SVD), the four candidate
decompositions are:
\[
  (\mat{R}, \vect{t}) \in \left\{
    (\mat{U}\mat{W}\mat{V}^\T,\, +\vect{u}_3),\;
    (\mat{U}\mat{W}\mat{V}^\T,\, -\vect{u}_3),\;
    (\mat{U}\mat{W}^\T\mat{V}^\T,\, +\vect{u}_3),\;
    (\mat{U}\mat{W}^\T\mat{V}^\T,\, -\vect{u}_3)
  \right\}
\]
where $\mat{W}=\begin{bmatrix}0&-1&0\\1&0&0\\0&0&1\end{bmatrix}$ and
$\vect{u}_3$ is the third column of $\mat{U}$.
The correct decomposition is chosen via the \emph{cheirality check}:
triangulated points must lie in front of both cameras ($Z>0$).
\end{mathbox}

\section{Kabsch Algorithm (SVD-Based Rotation)}

Used in \file{stereo\_vision.py} lines 188--209 to align two 3-D point sets.

\begin{theorem}[Kabsch Optimal Rotation]
Given two zero-centered point sets $\mat{P}, \mat{Q} \in \R^{n\times3}$, the
rotation $\mat{R}^*$ minimising $\sum_i \norm{\mat{R}\vect{p}_i - \vect{q}_i}^2$ is:
\[
  \mat{H} = \mat{P}^\T \mat{Q}, \quad \mat{H} = \mat{U}\mat{S}\mat{V}^\T \quad(\text{SVD})
\]
\[
  \mat{R}^* = \mat{V} \begin{bmatrix} 1 & 0 & 0 \\ 0 & 1 & 0 \\ 0 & 0 & \det(\mat{V}\mat{U}^\T) \end{bmatrix} \mat{U}^\T
\]
The $\det$ correction ensures $\det(\mat{R}^*)=+1$ (proper rotation).
\end{theorem}

\begin{lstlisting}
H = points_1_centered.T @ points_2_centered  # 3x3 cross-covariance
U, S, Vt = np.linalg.svd(H)
R = Vt.T @ U.T                               # initial estimate
if np.linalg.det(R) < 0:                     # reflection check
    Vt[-1, :] *= -1                           # flip last row
    R = Vt.T @ U.T                           # corrected rotation
t = centroid_2 - R @ centroid_1              # optimal translation
\end{lstlisting}

\section{Reinforcement Learning Foundations}

\subsection{Markov Decision Process}

\begin{definition}[MDP]
A \emph{Markov Decision Process} is a tuple $(\mathcal{S}, \mathcal{A}, P, R, \gamma)$:
\begin{itemize}
  \item $\mathcal{S}$: state space (localization uncertainty, energy need, battery mode)
  \item $\mathcal{A}$: action space $\{$\code{SKIP}, \code{CREATE\_KF}$\}$
  \item $P(s'|s,a)$: transition probability
  \item $R(s,a)$: reward function
  \item $\gamma \in [0,1)$: discount factor
\end{itemize}
\end{definition}

\subsection{Q-Learning Update Rule}

The tabular Q-learning update (used in \file{rl\_keyframe\_selector.py} line 258):
\[
  Q(s,a) \;\leftarrow\; Q(s,a) + \alpha \Bigl[ r + \gamma \max_{a'} Q(s',a') - Q(s,a) \Bigr]
\]
\begin{itemize}
  \item $\alpha$: learning rate (0.1 in code)
  \item $\gamma$: discount factor (0.95 in code)
  \item $r$: reward from \code{compute\_reward()}
  \item $\max_{a'} Q(s',a')$: bootstrapped future value from target state
\end{itemize}

\subsection{Deep Q-Network (Double DQN)}

Used in \file{dqn\_agent.py}. The Bellman target for Double DQN avoids
overestimation by decoupling action selection (policy net) from action evaluation
(target net):
\[
  y_i = r_i + (1 - d_i)\,\gamma\; Q_{\theta^-}\!\left(s'_i,\; \argmax_{a'} Q_\theta(s'_i, a')\right)
\]
\begin{itemize}
  \item $Q_\theta$: online (policy) network with parameters $\theta$
  \item $Q_{\theta^-}$: target network with lagged parameters $\theta^-$
  \item $d_i \in \{0,1\}$: done flag
\end{itemize}

Loss:
\[
  \mathcal{L}(\theta) = \frac{1}{|\mathcal{B}|} \sum_{i \in \mathcal{B}} \bigl(y_i - Q_\theta(s_i, a_i)\bigr)^2
\]

\subsection{Epsilon-Greedy Exploration}

\[
  a = \begin{cases}
    \text{random action} & \text{with probability } \varepsilon \\
    \argmax_a Q(s,a) & \text{with probability } 1-\varepsilon
  \end{cases}
\]

Linear decay used in the DQN agent:
\[
  \varepsilon(t) = (\varepsilon_\text{start} - \varepsilon_\text{end})\cdot
                  \left(1 - \frac{t}{T_\text{decay}}\right) + \varepsilon_\text{end},
  \quad t \le T_\text{decay}
\]

%=============================================================================
\chapter{Module: \texttt{camera\_stream\_handler.py}}
%=============================================================================

\section{Module Purpose}

This module provides two classes:
\begin{itemize}
  \item \code{CameraStreamHandler} --- spawns a background thread to continuously
        receive and decode MJPEG frames from a single IP camera URL.
  \item \code{StereoCameraManager} --- wraps two \code{CameraStreamHandler} instances
        and exposes a synchronised stereo frame pair.
\end{itemize}

\section{Line-by-Line Analysis: \texttt{CameraStreamHandler}}

\subsection{Imports and Header}

\begin{lstlisting}
"""
Camera Stream Handler Module            # L1-4: module docstring
Handles fetching and decoding MJPEG streams from multiple camera sources
"""

import cv2        # OpenCV: image I/O, decoding
import numpy as np  # NumPy: array operations
import threading  # Python threading for background stream
import time       # time.sleep(), time.time()
import requests   # HTTP streaming GET request
from collections import deque  # fixed-size circular buffer
from typing import Optional, Tuple  # type hints
\end{lstlisting}

\begin{notebox}[Why \texttt{requests} instead of OpenCV's VideoCapture?]
\code{cv2.VideoCapture(url)} is convenient but blocks the main thread if the
camera is unreachable and does not support graceful reconnection.
\code{requests.get(..., stream=True)} gives byte-level control, enabling custom
JPEG boundary detection and retry logic.
\end{notebox}

\subsection{Class \texttt{CameraStreamHandler}}

\begin{lstlisting}
class CameraStreamHandler:
    def __init__(self, camera_url: str, buffer_size: int = 2,
                 timeout: int = 5):
        self.camera_url = camera_url     # L28: store URL
        self.buffer_size = buffer_size   # L29: max frames in deque
        self.timeout = timeout           # L30: HTTP connect timeout (s)
        self.frame_buffer = deque(maxlen=buffer_size)  # L31: circular buffer
        self.latest_frame = None         # L32: most-recent decoded frame
        self.is_running = False          # L33: run/stop flag
        self.thread = None               # L34: streaming thread reference
        self.error_count = 0             # L35: consecutive error counter
        self.max_errors = 5              # L36: max retries before giving up
\end{lstlisting}

\textbf{Design choice -- \code{deque(maxlen=2)}:} A deque with \code{maxlen}
automatically evicts the oldest element when full, providing a lightweight
``always get the latest'' buffer without manual locking. For real-time SLAM it
is critical to process the \emph{most recent} frame, not a stale one.

\subsection{MJPEG Boundary Detection (Core Algorithm)}

\begin{lstlisting}
bytes_data = b''                           # L59: accumulator
for chunk in response.iter_content(chunk_size=8192):  # L60: 8KB chunks
    if not self.is_running:
        break
    bytes_data += chunk                    # L64: append

    a = bytes_data.find(b'\xff\xd8')       # L67: JPEG Start-of-Image (SOI)
    b = bytes_data.find(b'\xff\xd9')       # L68: JPEG End-of-Image (EOI)

    if a != -1 and b != -1:
        jpg_data = bytes_data[a:b+2]       # L71: slice one JPEG frame
        bytes_data = bytes_data[b+2:]      # L72: discard processed bytes

        frame = cv2.imdecode(
            np.frombuffer(jpg_data, dtype=np.uint8),
            cv2.IMREAD_COLOR                # L75: decode JPEG -> BGR array
        )
        if frame is not None:
            self.latest_frame = frame       # L77: thread-safe write
            self.frame_buffer.append(frame.copy())  # L78: buffered copy
            self.error_count = 0            # L79: reset error counter
\end{lstlisting}

\begin{mathbox}[MJPEG Frame Boundary]
JPEG files always start with the 2-byte \emph{Start of Image} (SOI) marker
\texttt{0xFF 0xD8} and end with the \emph{End of Image} (EOI) marker
\texttt{0xFF 0xD9}. An MJPEG stream concatenates multiple JPEGs in sequence.
Scanning for these delimiters extracts individual frames even from partial
network chunks.
\end{mathbox}

\subsection{Automatic Reconnection}

\begin{lstlisting}
except Exception as e:
    self.error_count += 1
    print(f"Stream error ({self.camera_url}): {str(e)}")
    if self.error_count < self.max_errors:
        time.sleep(2)                    # L85: back-off delay
        if self.is_running:
            self._stream_worker()        # L87: recursive retry
\end{lstlisting}

\noindent The retry logic uses \emph{recursive self-call} with a 2-second
back-off. This is simple but note that Python has a default recursion limit of
1000; with 5 max errors and 2-second sleep the limit is never approached in
practice.

\section{Class \texttt{StereoCameraManager}}

\begin{lstlisting}
class StereoCameraManager:
    def __init__(self, left_camera_url: str, right_camera_url: str):
        self.left_camera = CameraStreamHandler(left_camera_url)   # L115
        self.right_camera = CameraStreamHandler(right_camera_url) # L116
\end{lstlisting}

Synchronisation strategy (lines 125--133):
\begin{lstlisting}
timeout = 10
start_time = time.time()
while time.time() - start_time < timeout:
    if (self.left_camera.is_connected() and
            self.right_camera.is_connected()):
        return True
    time.sleep(0.5)
\end{lstlisting}

\begin{warnbox}[Temporal Synchronisation Limitation]
The two cameras are \emph{not} hardware-synchronised. Frames are fetched
independently; the pairing is ``latest left + latest right''. This can
introduce a temporal offset $\Delta t \approx 0$--$100$\,ms at 10\,fps.
For a camera travelling at $0.5$\,m/s, $\Delta t = 50$\,ms yields a
position error of $\approx 2.5$\,cm --- acceptable for indoor SLAM but
not for high-speed outdoor operation.
\end{warnbox}

%=============================================================================
\chapter{Module: \texttt{stereo\_vision.py}}
%=============================================================================

\section{Module Purpose}

Implements the \code{StereoCalibratedPipeline} class, which:
\begin{enumerate}
  \item Extracts SIFT keypoints from a frame.
  \item Matches features between frames using Lowe's ratio test.
  \item Computes a dense disparity map using \code{StereoBM}.
  \item Triangulates 3-D points from the disparity map.
  \item Estimates camera motion between two point sets (Kabsch SVD).
\end{enumerate}

\section{SIFT Feature Extraction}

\begin{lstlisting}
self.sift = cv2.SIFT_create()    # L27: create SIFT detector
\end{lstlisting}

\begin{defbox}[SIFT: Scale-Invariant Feature Transform (Lowe, 2004)]
SIFT detects keypoints at extrema of a \emph{Difference-of-Gaussian} (DoG)
scale-space. For scale $\sigma$:
\[
  D(x,y,\sigma) = G(x,y,k\sigma)*I - G(x,y,\sigma)*I
\]
where $G$ is a Gaussian kernel. Each keypoint is assigned a dominant gradient
orientation, making the $128$-dimensional descriptor rotation-invariant.
\end{defbox}

\subsection{Lowe's Ratio Test}

\begin{lstlisting}
matches = self.bf_matcher.knnMatch(desc1, desc2, k=2)  # L81
good_matches = []
for match_pair in matches:
    if len(match_pair) == 2:
        m, n = match_pair        # nearest and second-nearest
        if m.distance < ratio_threshold * n.distance:  # L88: ratio test
            good_matches.append((m.queryIdx, m.trainIdx))
\end{lstlisting}

\begin{mathbox}[Lowe's Ratio Test]
For each descriptor $\vect{d}_i$ in image 1, let $\vect{d}^1_j$ and
$\vect{d}^2_j$ be the closest and second-closest descriptors in image 2.
A match is accepted iff:
\[
  \frac{\norm{\vect{d}_i - \vect{d}^1_j}}{\norm{\vect{d}_i - \vect{d}^2_j}} < \tau, \qquad \tau = 0.7
\]
This rejects ambiguous matches where two descriptors are nearly equidistant,
reducing false matches from $\sim$46\% to $\sim$10\% (Lowe, 2004).
\end{mathbox}

\section{Dense Disparity with \texttt{StereoBM}}

\begin{lstlisting}
self.stereo_matcher = cv2.StereoBM_create(
    numDisparities=128,  # L34: search range [0, 128] pixels
    blockSize=15         # L35: 15x15 correlation window
)
\end{lstlisting}

\noindent \code{StereoBM} implements \emph{Semi-Global Block Matching} (SGBM).
For each pixel $(u,v)$ it finds the horizontal shift $d$ that minimises the
Sum of Absolute Differences (SAD) within a block:

\[
  \hat{d}(u,v) = \argmin_{d} \sum_{(i,j)\in W} \abs{I_L(u+i,\; v+j) - I_R(u+i-d,\; v+j)}
\]

\section{3-D Triangulation}

\begin{lstlisting}
# Lines 155-168
depth = depth_map[y_l, x_l]
if depth > 0:
    cx, cy = self.frame_width / 2, self.frame_height / 2
    X = (x_l - cx) * depth / self.focal_length
    Y = (y_l - cy) * depth / self.focal_length
    Z = depth
    points_3d.append([X, Y, Z])
\end{lstlisting}

Each pixel $(u,v)$ with depth $Z$ back-projects to 3-D using the inverse
pinhole model:
\[
  \begin{bmatrix}X\\Y\\Z\end{bmatrix} =
  Z \cdot \mat{K}^{-1} \begin{bmatrix}u\\v\\1\end{bmatrix} =
  \begin{bmatrix}\frac{(u-c_x)Z}{f} \\ \frac{(v-c_y)Z}{f} \\ Z\end{bmatrix}
\]

%=============================================================================
\chapter{Module: \texttt{advanced\_features.py}}
%=============================================================================

\section{Overview}

This module contains six classes:
\begin{enumerate}
  \item \code{AdvancedFeatureExtractor} --- unified SIFT/ORB/AKAZE/hybrid extractor with CLAHE preprocessing
  \item \code{RobustStereoMatcher} --- BF or FLANN matcher with ratio test
  \item \code{EpipolarGeometry} --- fundamental matrix and epipolar filtering
  \item \code{PointCloudFilter} --- outlier removal, voxel downsampling, deduplication
  \item \code{LoopClosureDetector} --- geometric loop closure via Euclidean distance
  \item \code{BundleAdjustment} --- simplified iterative pose refinement
  \item \code{CloudExporter} --- PLY and JSON export utilities
\end{enumerate}

\section{CLAHE Preprocessing}

\begin{lstlisting}
clahe = cv2.createCLAHE(clipLimit=2.0, tileGridSize=(8, 8))  # L62
gray = clahe.apply(gray)                                       # L63
\end{lstlisting}

\begin{defbox}[CLAHE: Contrast Limited Adaptive Histogram Equalization]
Standard histogram equalization (HE) can over-amplify noise. CLAHE divides
the image into $8\times8$ tiles and applies HE independently in each tile,
clipping the histogram at \emph{clipLimit} to limit amplification. Bilinear
interpolation then merges tile boundaries. Result: locally enhanced contrast
without noise explosion, improving feature detection in dark or uneven
illumination.
\end{defbox}

\section{Feature Quality Score}

\begin{lstlisting}
# _calculate_quality_score (lines 93-111)
num_score = min(len(keypoints) / self.max_features, 1.0)    # L99

positions = np.array([kp.pt for kp in keypoints])
hist, _ = np.histogram(positions[:, 0], bins=10)            # L105
entropy = -np.sum((hist/hist.sum())*np.log(hist/hist.sum()+1e-10))  # L106
dist_score = entropy / np.log(10)                            # L107

return 0.7 * num_score + 0.3 * dist_score                   # L111
\end{lstlisting}

\begin{mathbox}[Spatial Entropy]
Spatial entropy measures how uniformly keypoints are distributed across the
image width. Let $h_k$ be the count in bin $k$ (10 bins along $x$-axis):
\[
  H = -\sum_{k=1}^{10} p_k \ln p_k, \quad p_k = \frac{h_k}{\sum_j h_j}
\]
Maximum entropy $H_{\max} = \ln 10$ occurs when keypoints are perfectly
uniform. The quality component is $H / H_{\max} \in [0,1]$.
A high-quality frame has many ($70\%$ weight) and spatially distributed
($30\%$ weight) features.
\end{mathbox}

\section{Epipolar Constraint Filtering}

\begin{lstlisting}
# filter_by_epipolar_constraint (lines 241-266)
points1_h = np.hstack([points1, np.ones((len(points1), 1))])  # homogeneous
lines2 = F @ points1_h.T                      # epipolar lines in img 2
distances = np.abs(np.sum(points2 * lines2[:2].T, axis=1)) / \
            np.sqrt(lines2[0]**2 + lines2[1]**2)   # point-to-line dist
return distances < threshold
\end{lstlisting}

\begin{mathbox}[Epipolar Distance]
A point $\vect{x}_1 = [u_1,v_1,1]^\T$ maps to epipolar line $\vect{l}_2 = \mat{F}\vect{x}_1$
in image 2. The algebraic line $\vect{l}_2 = [a,b,c]^\T$ is defined by
$au+bv+c=0$. The distance from point $\vect{x}_2=[u_2,v_2,1]^\T$ to this
line is:
\[
  d(\vect{x}_2, \vect{l}_2) = \frac{\abs{a u_2 + b v_2 + c}}{\sqrt{a^2+b^2}}
\]
Matches with $d > 1.0$ pixel are rejected as geometric outliers.
\end{mathbox}

\section{Voxel Grid Downsampling}

\begin{lstlisting}
indices = np.floor(points / voxel_size).astype(int)   # L320
_, unique_idx = np.unique(indices, axis=0,
                          return_index=True)            # L323
return points[unique_idx], colors[unique_idx]
\end{lstlisting}

\begin{defbox}[Voxel Grid Downsampling]
The 3-D space is partitioned into cubic voxels of side $v$ metres. Each point
is mapped to the voxel it falls in via $\lfloor \vect{p}/v \rfloor$.
Only the first point per unique voxel index is kept. This reduces $N$ points
to at most $\lceil \text{vol}/v^3 \rceil$ points while preserving spatial
structure, running in $\mathcal{O}(N\log N)$ time.
\end{defbox}

%=============================================================================
\chapter{Module: \texttt{visual\_odometry.py}}
%=============================================================================

\section{Purpose}

The \code{VisualOdometryEngine} tracks the cumulative 6-DoF camera pose by:
\begin{enumerate}
  \item \emph{Frame 0}: initialise 3-D point cloud from left-right stereo pair.
  \item \emph{Frame $t>0$}: match features between current and previous left frames,
        estimate $(\mat{R},\vect{t})$ via essential matrix, accumulate pose.
\end{enumerate}

\section{Class Initialization}

\begin{lstlisting}
self.stereo_pipeline = StereoCalibratedPipeline()   # L24: reuse stereo vision
self.camera_pose = np.eye(4)   # L27: identity = world origin = camera start
self.trajectory = deque(maxlen=1000)   # L30
self.point_cloud = np.array([]).reshape(0, 3)   # L34: empty Nx3 array
self.prev_keypoints = None   # L38: will store prior frame's keypoints
\end{lstlisting}

\noindent Starting from the identity matrix means the world coordinate system
coincides with the camera at time 0.

\section{Pose Update Derivation}

Given relative rotation $\mat{R}$ and translation $\vect{t}$ between frames
$t-1$ and $t$:

\begin{lstlisting}
# _update_camera_pose (lines 259-265)
T = np.eye(4)
T[:3, :3] = R       # L261: fill rotation block
T[:3, 3] = t.flatten()  # L262: fill translation column

self.camera_pose = self.camera_pose @ T   # L265: right-multiply to accumulate
\end{lstlisting}

Mathematically, if $\mat{T}_k$ is the pose at frame $k$ (world-to-camera), then:
\[
  \mat{T}_{k+1} = \mat{T}_k \cdot \mat{T}_{\Delta}, \quad
  \mat{T}_\Delta = \begin{bmatrix}\mat{R} & \vect{t}\\\vect{0}^\T & 1\end{bmatrix}
\]

The camera centre in world coordinates extracted as:
\begin{lstlisting}
camera_center = pose[:3, 3]   # L274
\end{lstlisting}
This is the translation column of the accumulated pose, i.e.,
$\vect{c}_w = \vect{t}_\text{acc}$.

\section{Point Cloud World Transformation}

\begin{lstlisting}
# _transform_to_world (lines 279-296)
R = self.camera_pose[:3, :3]
t = self.camera_pose[:3, 3]
points_world = (R @ points_3d.T).T + t    # L296
\end{lstlisting}

\[
  \vect{p}_w = \mat{R}\,\vect{p}_c + \vect{t}
\]

This is a \emph{vectorised} operation: \code{points\_3d.T} is $3\times N$,
so \code{R @ points\_3d.T} is $3\times N$, then \code{.T} gives $N\times3$,
and \code{+ t} broadcasts the $3$-vector to every row.

%=============================================================================
\chapter{Module: \texttt{stereo\_slam.py}}
%=============================================================================

\section{Purpose and Architecture}

\code{StereoSLAMSystem} is the central SLAM engine. It maintains:
\begin{itemize}
  \item A deque of \code{KeyFrame} objects (circular buffer, max 100).
  \item A growing 3-D point cloud and associated colors.
  \item A trajectory (deque of 3-D camera positions).
  \item A loop-closure graph.
\end{itemize}

\section{KeyFrame Class}

\begin{lstlisting}
class KeyFrame:
    def __init__(self, frame_id, timestamp, pose, left_frame,
                 right_frame, keypoints, descriptors):
        self.frame_id = frame_id          # L34: unique integer ID
        self.timestamp = timestamp        # L35: Unix time of capture
        self.pose = pose.copy()           # L36: 4x4 SE(3) matrix
        self.left_frame = left_frame.copy()   # L37: raw BGR image
        self.right_frame = right_frame.copy() # L38: raw BGR image
        self.keypoints = keypoints        # L39: list of cv2.KeyPoint
        self.descriptors = descriptors.copy() # L40: NxD numpy array
        self.point_ids = {}               # L41: keypoint_idx -> point3D_id
\end{lstlisting}

\noindent All arrays are \emph{deep-copied} to avoid aliasing bugs when frames
are recycled.

\begin{lstlisting}
def get_camera_center(self) -> np.ndarray:
    return self.pose[:3, 3]    # L45: extract translation = camera position
\end{lstlisting}

\section{Keyframe Decision}

\begin{lstlisting}
def _should_create_keyframe(self, current_pose):
    if self.last_keyframe is None:
        return True   # L191: always accept first frame
    displacement = np.linalg.norm(
        current_pose[:3, 3] - self.last_pose[:3, 3])  # L194
    return displacement > self.min_displacement        # L196
\end{lstlisting}

\begin{mathbox}[Displacement Criterion]
The keyframe creation criterion is a simple Euclidean distance threshold in
translation space:
\[
  d = \norm{\vect{t}_{k} - \vect{t}_{k-1}}_2 > d_{\min}
\]
where $d_{\min} = 0.05$\,m by default. This ensures keyframes capture
sufficient parallax for reliable triangulation. The RL agent and battery manager
\emph{adaptively scale} $d_{\min}$ to trade off accuracy vs.\ energy.
\end{mathbox}

\section{Loop Closure Detection}

\begin{lstlisting}
# _detect_loop_closure (lines 266-302)
early_keyframes = list(self.keyframes)[:-10]   # L277: exclude last 10
for old_keyframe in early_keyframes:
    matches = self.stereo_matcher.match_temporal(
        old_keyframe.descriptors, keyframe.descriptors  # L284-286
    )
    match_score = len(matches) / (len(old_keyframe.keypoints)
                                  + len(keyframe.keypoints))   # L290
    if match_score > self.loop_closure_threshold:   # L297: 0.75
        self.loop_closures.append(...)
\end{lstlisting}

\begin{defbox}[Loop Closure via Descriptor Matching]
Loop closure detects when the camera revisits a previously seen location.
The match score is defined as:
\[
  s = \frac{|\mathcal{M}|}{n_{\text{old}} + n_{\text{current}}}
\]
where $|\mathcal{M}|$ is the number of successful descriptor matches and
$n_{\text{old}}, n_{\text{current}}$ are the total keypoints. A high score
indicates visual similarity consistent with loop closure.
\end{defbox}

\section{Point Cloud Integration}

\begin{lstlisting}
# _add_points_from_stereo (lines 198-237)
points_3d = self.visual_odometry.stereo_pipeline
            .compute_3d_points_from_matches(...)   # L212-214

points_world = self._transform_to_world(points_3d, keyframe.pose)  # L218
colors_sampled = self._sample_colors(left_frame, left_kp, matches) # L221

self.point_cloud = np.vstack([self.point_cloud, points_world])      # L224
self.point_cloud_colors = np.vstack([self.point_cloud_colors,
                                     colors_sampled])                # L225
\end{lstlisting}

\noindent \code{np.vstack} concatenates along the row axis, growing the point
cloud incrementally. In memory, a new array is allocated each time; for
very large maps consider pre-allocated arrays with index tracking.

\section{Bundle Adjustment (Simplified)}

Full bundle adjustment minimises the reprojection error over all keyframes
and 3-D points simultaneously:
\[
  \min_{\{\mat{T}_k\},\{\vect{P}_j\}} \sum_{k,j} \norm{\pi(\mat{T}_k, \vect{P}_j) - \vect{x}_{kj}}^2
\]
where $\pi$ is the projection function and $\vect{x}_{kj}$ is the observed
2-D feature location. The current implementation (\code{optimize\_bundle},
lines 304--323) provides a placeholder that iterates over the last 5 keyframes
without actual optimisation --- a future extension point for a full Gauss-Newton
or g2o solver.

%=============================================================================
\chapter{Module: \texttt{rl\_keyframe\_selector.py}}
%=============================================================================

\section{Purpose}

This is the \emph{Q-learning brain} of the system. It defines:
\begin{enumerate}
  \item \code{KeyframeAction} --- a 2-action enum (\code{SKIP\_FRAME}, \code{CREATE\_KEYFRAME}).
  \item \code{LocalizationUncertaintyMetrics} --- computes a scalar uncertainty $u \in [0,1]$ from the SLAM state.
  \item \code{EnergyModel} --- returns a scalar energy-need $e \in [0,1]$ from the battery mode.
  \item \code{RLKeyframeSelector} --- tabular Q-learning agent with discretised state space.
  \item \code{RL\_SLAM\_Integration} --- thin integration layer connecting RL decisions to the SLAM main loop.
\end{enumerate}

\section{State Space Design}

\begin{lstlisting}
def discretize_state(self, uncertainty: float,
                     energy_need: float) -> Tuple[int, int]:
    uncertainty_bin = min(int(uncertainty * self.state_bins),
                         self.state_bins - 1)   # L154
    energy_bin = min(int(energy_need * self.state_bins),
                    self.state_bins - 1)        # L155
    return uncertainty_bin, energy_bin
\end{lstlisting}

With \code{state\_bins=10} and 4 battery modes, the state space has
$10 \times 10 \times 4 = 400$ discrete states. The Q-table maps each state
to 2 action values: $(Q_{\text{skip}}, Q_{\text{create}})$.

\section{Uncertainty Metric Computation}

\begin{lstlisting}
# update_from_slam (lines 33-76)
# Component 1: position variance of last 5 trajectory points
recent_positions = trajectory[-5:]
self.position_variance = float(np.var(recent_positions))   # L43

# Component 2: tracking error from point density
points_per_kf = total_points / max(total_keyframes, 1)    # L51
tracking_error = 1.0 - (points_per_kf / 100.0) if < 100  # L55

# Component 3: feature match confidence (loop closure rate)
feature_match_confidence = max(0, 1.0-(loop_closures/...)) # L64

# Combined uncertainty
uncertainty = (
    min(self.position_variance, 1.0) * 0.3 +
    self.tracking_error * 0.4 +
    (1.0 - self.feature_match_confidence) * 0.3   # L67-71
)
\end{lstlisting}

\begin{mathbox}[Uncertainty Aggregation]
\[
  u = 0.3 \min(\sigma^2_\text{pos}, 1) + 0.4\, e_\text{track} + 0.3\,(1 - c_\text{feat})
\]
where:
\begin{itemize}
  \item $\sigma^2_\text{pos}$: variance of the last 5 camera positions
        (high if camera is drifting or estimations are noisy)
  \item $e_\text{track} = 1 - \min(\bar{n}_{\text{pts}}/100,\, 1)$: normalised
        inverse point density; sparse maps indicate poor tracking
  \item $c_\text{feat}$: feature confidence; decreases as loop-closure rate rises
        (many loop closures suggest uncorrected drift)
\end{itemize}
\end{mathbox}

\section{Reward Function}

\begin{lstlisting}
# compute_reward (lines 191-236)
uncertainty_improvement = max(0, prev_uncertainty - uncertainty)
localization_reward = uncertainty_improvement * 10.0   # L211

if action == CREATE_KEYFRAME:
    energy_penalty = -energy_need * 5.0               # L215
    if uncertainty > 0.5:
        localization_reward += 5.0    # L219: needed KF
    else:
        localization_reward -= 3.0    # L221: wasted KF
    reward = localization_reward + energy_penalty

else:  # SKIP_FRAME
    energy_reward = energy_need * 3.0    # L227
    if uncertainty > 0.6:
        energy_reward -= 5.0             # L231: missed needed KF
    reward = energy_reward

return np.clip(reward, -10.0, 10.0)     # L236
\end{lstlisting}

\begin{mathbox}[Reward Decomposition]
\[
  r = \underbrace{10 \max(0, u_{t-1}-u_t)}_{\text{localisation gain}}
    + \begin{cases}
        -5\,e + (\pm 5) & \text{CREATE\_KF} \\
        +3\,e - 5\cdot\mathbb{1}[u>0.6] & \text{SKIP}
      \end{cases}
\]
The reward is \emph{dense}: every step provides a signal, making learning
faster than sparse rewards (e.g., only rewarding at episode end).
\end{mathbox}

\section{Q-Learning Update}

\begin{lstlisting}
# update_q_value (lines 238-261)
current_q = self.q_table[state_key][action]   # L250

if done:
    max_next_q = 0.0
else:
    max_next_q = max(self.q_table[next_state_key].values())  # L255

# Q-learning Bellman update
new_q = current_q + self.learning_rate * (
    reward + self.discount_factor * max_next_q - current_q)  # L258
self.q_table[state_key][action] = new_q   # L259
\end{lstlisting}

This is the \textbf{exact Q-learning update}:
\[
  Q(s,a) \xleftarrow{\alpha} Q(s,a) + \alpha\bigl[r + \gamma \max_{a'}Q(s',a') - Q(s,a)\bigr]
\]
The term $\bigl[r + \gamma \max_{a'}Q(s',a') - Q(s,a)\bigr]$ is called the
\emph{temporal-difference (TD) error} $\delta$.

%=============================================================================
\chapter{Module: \texttt{dqn\_agent.py}}
%=============================================================================

\section{Purpose}

Implements a \textbf{Deep Q-Network (DQN)} with Double-DQN correction and
experience replay for the 3-action keyframe selection problem:
\begin{itemize}
  \item Action 0: \code{SKIP} (cheapest; skip frame entirely)
  \item Action 1: \code{SAVE\_NORMAL} (full keyframe)
  \item Action 2: \code{SAVE\_LIGHT} (pose-only keyframe)
\end{itemize}

\section{Network Architecture}

\begin{lstlisting}
class DQNNetwork(nn.Module):
    def __init__(self, state_dim=7, action_dim=3, hidden_dim=128):
        super(DQNNetwork, self).__init__()   # L34
        self.fc1 = nn.Linear(state_dim, hidden_dim)    # L36: 7 -> 128
        self.fc2 = nn.Linear(hidden_dim, hidden_dim)   # L37: 128 -> 128
        self.fc3 = nn.Linear(hidden_dim, action_dim)   # L38: 128 -> 3
        self.relu = nn.ReLU()                           # L40

    def forward(self, state):
        x = self.relu(self.fc1(state))   # L44
        x = self.relu(self.fc2(x))       # L45
        q_values = self.fc3(x)           # L46: no activation (regression)
        return q_values
\end{lstlisting}

\begin{mathbox}[Network Computation]
For state $\vect{s} \in \R^7$:
\[
  \vect{h}_1 = \mathrm{ReLU}(\mat{W}_1 \vect{s} + \vect{b}_1), \quad \mat{W}_1 \in \R^{128\times7}
\]
\[
  \vect{h}_2 = \mathrm{ReLU}(\mat{W}_2 \vect{h}_1 + \vect{b}_2), \quad \mat{W}_2 \in \R^{128\times128}
\]
\[
  \vect{q} = \mat{W}_3 \vect{h}_2 + \vect{b}_3, \quad \mat{W}_3 \in \R^{3\times128}
\]
Total parameters: $7\cdot128 + 128 + 128\cdot128 + 128 + 128\cdot3 + 3 = 18\,051$.
\end{mathbox}

\section{Experience Replay Buffer}

\begin{lstlisting}
class ReplayBuffer:
    def __init__(self, capacity: int = 100000):
        self.buffer = deque(maxlen=capacity)   # L60: circular buffer

    def push(self, state, action, reward, next_state, done):
        self.buffer.append((state, action, reward, next_state, done))  # L66

    def sample(self, batch_size: int) -> Tuple:
        batch = random.sample(self.buffer, batch_size)  # L73: iid sampling
        states = np.array([x[0] for x in batch])        # L75
        actions = np.array([x[1] for x in batch])       # L76
        rewards = np.array([x[2] for x in batch])       # L77
        next_states = np.array([x[3] for x in batch])   # L78
        dones = np.array([x[4] for x in batch])         # L79
        return states, actions, rewards, next_states, dones
\end{lstlisting}

\begin{notebox}[Why Experience Replay?]
Naive online learning with consecutive frames violates the i.i.d.\ assumption
of stochastic gradient descent because consecutive SLAM frames are highly
correlated. Replaying random past experiences breaks this correlation,
substantially improving training stability and sample efficiency.
\end{notebox}

\section{Double DQN Training Step}

\begin{lstlisting}
# train_step (lines 180-240)

# Current Q-values for taken actions
q_values = self.policy_net(states)             # L208: [B, 3]
q_values = q_values.gather(1, actions.unsqueeze(1)).squeeze(1)  # L209

# Double DQN target (no gradient)
with torch.no_grad():
    next_q_values = self.target_net(next_states)  # L213: target net
    next_actions = self.policy_net(next_states).argmax(1)    # L214: policy selects
    next_q_values = next_q_values.gather(1,
                    next_actions.unsqueeze(1)).squeeze(1)    # L215: target evaluates
    target_q_values = (rewards
                       + (1 - dones) * self.gamma * next_q_values)  # L216

target_q_values = torch.clamp(target_q_values, -50, 50)  # L219: stability

loss = nn.MSELoss()(q_values, target_q_values)  # L222

self.optimizer.zero_grad()                      # L225
loss.backward()                                 # L226: backprop
torch.nn.utils.clip_grad_norm_(
    self.policy_net.parameters(), 1.0)          # L227: gradient clipping
self.optimizer.step()                           # L228
\end{lstlisting}

\begin{mathbox}[Double DQN vs. Standard DQN]
\textbf{Standard DQN} (Mnih et al., 2015):
\[
  y^{\text{DQN}} = r + \gamma \max_{a'} Q_{\theta^-}(s', a')
\]
\textbf{Double DQN} (van Hasselt et al., 2016):
\[
  y^{\text{DDQN}} = r + \gamma Q_{\theta^-}\!\bigl(s',\; \argmax_{a'}Q_\theta(s',a')\bigr)
\]
Standard DQN overestimates Q-values because it uses the same network to select
\emph{and} evaluate the greedy action. Double DQN separates these roles:
the online network selects; the target network evaluates. This reduces
overestimation bias by $\sim20$--$40\%$ on typical RL benchmarks.
\end{mathbox}

\subsection{Gradient Clipping}

\begin{lstlisting}
torch.nn.utils.clip_grad_norm_(self.policy_net.parameters(), 1.0)
\end{lstlisting}

Clips the global gradient norm to 1.0:
\[
  \text{if } \norm{\nabla_\theta \mathcal{L}}_2 > 1 \implies
  \nabla_\theta \leftarrow \frac{\nabla_\theta}{\norm{\nabla_\theta}_2}
\]
This prevents exploding gradients, a common failure mode in RL training where
high-variance rewards produce large, destabilising parameter updates.

\subsection{Target Network Update}

\begin{lstlisting}
if self.training_steps % self.target_update_freq == 0:  # L234: every 1000 steps
    self.target_net.load_state_dict(self.policy_net.state_dict())  # L235
\end{lstlisting}

The target network $\theta^-$ is a lagged copy of $\theta$, updated every
$T_{\text{target}}=1000$ steps. This stabilises training by providing
non-moving regression targets.

\subsection{Epsilon Decay}

\begin{lstlisting}
# _update_epsilon (lines 242-249)
if self.step_count < self.epsilon_decay_steps:
    self.epsilon = (self.epsilon_start - self.epsilon_end) * \
                   (1 - self.step_count / self.epsilon_decay_steps) \
                   + self.epsilon_end   # L245-246
else:
    self.epsilon = self.epsilon_end     # L248: plateau at min
self.step_count += 1
\end{lstlisting}

Linear annealing from $\varepsilon_0 = 0.3$ to $\varepsilon_\infty = 0.05$
over $T = 100\,000$ steps:
\[
  \varepsilon(t) = (\varepsilon_0 - \varepsilon_\infty)\left(1 - \frac{t}{T}\right) + \varepsilon_\infty
\]

%=============================================================================
\chapter{Module: \texttt{energy\_battery\_model.py}}
%=============================================================================

\section{Energy Cost Table}

\begin{lstlisting}
ENERGY_COSTS = {
    'skip_frame':        0.02,   # minimal tracking only
    'save_light':        0.20,   # pose + reduced descriptors
    'save_normal':       0.60,   # full descriptor set
    'save_high':         1.00,   # full + extra processing
    'bundle_adjustment': 0.30,
    'loop_closure_check':0.15,
    'map_transmission':  0.40
}
\end{lstlisting}

These normalised costs (relative to the worst-case ``save\_high'') are used in
reward computation. They are engineering estimates; in production they
should be calibrated with actual power measurements via a hardware power meter.

\section{Battery-Aware Energy Weight Multiplier}

\begin{lstlisting}
def get_energy_weight_multiplier(self, k_batt: float = 1.5) -> float:
    norm_battery = self.get_normalized_battery()      # b in [0,1]
    return 1.0 + k_batt * (1.0 - norm_battery)       # L192
\end{lstlisting}

\begin{mathbox}[Energy Weight Multiplier]
\[
  w_e(b) = 1 + k_{\text{batt}}\,(1 - b), \quad b \in [0,1]
\]
At full battery ($b=1$): $w_e = 1$ (energy cost at face value).\\
At empty battery ($b=0$): $w_e = 1 + 1.5 = 2.5$ (costs 2.5$\times$ more).\\
This continuously scales the penalty for expensive actions as battery depletes.
\end{mathbox}

\section{Dense Reward Calculator}

\begin{lstlisting}
# compute_reward (lines 263-325)
delta_uncertainty = np.clip(prev_u - curr_u, -MAX_LOSS, MAX_GAIN)  # L287
localization_reward = delta_uncertainty * W_BENEFIT                 # L289: *8

energy_cost = self.energy_model.get_action_cost(action)             # L292
energy_weight = W_ENERGY * battery_monitor.get_energy_weight_multiplier()  # L296
energy_penalty = -energy_weight * energy_cost                       # L300

if not tracking_ok:
    tracking_penalty = -W_TRACKING_LOSS   # L304: -40 for lost tracking

step_penalty = -W_STEP   # L309: -0.01 (dithering discouragement)

# Bonus for successful skip; penalty for wasted save
bonus = BONUS_SKIP if skip_ok else -PENALTY_WASTE_if_wasted   # L313-317

reward = localization_r + energy_penalty + tracking_penalty + step_penalty + bonus
reward = np.clip(reward, -50.0, 50.0)
\end{lstlisting}

Complete reward formula:
\[
  r = 8\,\Delta u_{\text{clip}} - W_E\,w_b\,c_a - 40\,\mathbb{1}[\neg\text{track}] - 0.01 + \text{bonus}
\]
where $\Delta u_{\text{clip}} = \text{clip}(u_{t-1}-u_t, -0.2, 0.2)$ and
$c_a$ is the normalised action energy cost.

%=============================================================================
\chapter{Module: \texttt{degradation\_policy.py}}
%=============================================================================

\section{4-Level Degradation Policy}

\begin{table}[H]
\centering
\caption{SLAM degradation levels vs.\ battery range}\label{tab:degrad}
\begin{tabular}{clcccccc}
\toprule
\textbf{Level} & \textbf{Name} & \textbf{Batt.\%} & \textbf{Desc.} & \textbf{KF freq.} & \textbf{LC freq.} & \textbf{Max desc.} & \textbf{BA iter.} \\
\midrule
0 & FULL FIDELITY & 60--100 & SIFT & 100\% & 100\% & 2000 & 5 \\
1 & MILD & 40--60 & SIFT & 90\% & 90\% & 1500 & 3 \\
2 & MODERATE & 20--40 & ORB & 70\% & 50\% & 800 & 1 \\
3 & AGGRESSIVE & 0--20 & FAST & 40\% & 10\% & 300 & 0 \\
\bottomrule
\end{tabular}
\end{table}

\noindent Level transitions are determined by:
\begin{lstlisting}
def _get_level_from_battery(self, battery_pct: float) -> int:
    if battery_pct >= 60: return 0
    elif battery_pct >= 40: return 1
    elif battery_pct >= 20: return 2
    else: return 3
\end{lstlisting}

\section{Stochastic Keyframe Gating}

\begin{lstlisting}
def should_create_keyframe(self, random_value: float = None) -> bool:
    if random_value is None:
        random_value = np.random.random()
    threshold = self.configuration['keyframe_frequency']
    return random_value < threshold   # L175
\end{lstlisting}

At level 3 (\code{keyframe\_frequency=0.4}), each frame independently has a
$40\%$ probability of being promoted to a keyframe. This is a simple
\emph{Bernoulli trial}: $X \sim \text{Bernoulli}(p=0.4)$.

\section{Descriptor Subsampling}

\begin{lstlisting}
descriptor_variance = np.var(descriptors, axis=1)         # L250: per-descriptor
top_indices = np.argsort(descriptor_variance)[-keep_count:]  # L251: highest var
return descriptors[top_indices]                              # L253
\end{lstlisting}

\begin{notebox}[Information-Theoretic Justification]
High-variance descriptors encode more discriminative information. Retaining
the top-$k\%$ by descriptor-dimension variance is a heuristic approximation
to maximising the entropy of the retained descriptor set, preserving the most
unique feature information given a fixed budget.
\end{notebox}

%=============================================================================
\chapter{Module: \texttt{integrated\_slam.py} (Root Entry Point)}
%=============================================================================

\section{Purpose}

The root \file{integrated\_slam.py} is the entry point for the \emph{non-RL}
version of the system. It wires together:
\begin{enumerate}
  \item \code{BatteryManager} --- monitors system battery and enforces hard frame-skipping rules.
  \item \code{IntegratedSLAMVisualizer} --- real-time 3-D matplotlib animation.
  \item \code{IntegratedStereoSLAMSystem} --- main run loop.
\end{enumerate}

\section{BatteryManager}

\begin{lstlisting}
BATTERY_CRITICAL = 5   # L31: < 5%  -> process 1/8 frames
BATTERY_LOW      = 15  # L32: < 15% -> process 1/4 frames
BATTERY_NORMAL   = 50  # L33: < 50% -> process 1/2 frames
BATTERY_GOOD     = 100 # L34: >= 50% -> process all frames

def should_skip_frame(self, frame_count: int) -> bool:
    mode = self.get_mode()
    if mode == 'critical':
        return frame_count % 8 != 0   # L72-73: 1 of 8
    elif mode == 'low':
        return frame_count % 4 != 0   # L75-76: 1 of 4
    elif mode == 'moderate':
        return frame_count % 2 != 0   # L78-79: 1 of 2
    else:
        return False                   # L81-82: all frames
\end{lstlisting}

This is a \emph{deterministic} gating policy (as opposed to the stochastic
Bernoulli policy in \code{degradation\_policy.py}) based on the frame counter
modulo.

\subsection{Keyframe Displacement Multiplier}

\begin{lstlisting}
def get_keyframe_threshold_multiplier(self) -> float:
    mode = self.get_mode()
    if mode == 'critical': return 3.0   # L89
    elif mode == 'low':    return 2.0   # L91
    elif mode == 'moderate':return 1.5  # L93
    else:                  return 1.0   # L95
\end{lstlisting}

Adaptive threshold:
\[
  d_{\min}^{\text{eff}} = 0.05 \times m_{\text{batt}}
\]
At \code{critical} mode, the minimum displacement becomes 0.15\,m (3$\times$ larger),
so keyframes are only created when the camera has moved significantly, drastically
reducing the keyframe creation rate.

\section{3-D Visualizer: FuncAnimation}

\begin{lstlisting}
ani = FuncAnimation(self.visualizer.fig, update_plot,
                    interval=300,              # L459: 300ms between frames
                    blit=True,                 # only re-draw changed artists
                    cache_frame_data=False)    # no caching (live data)
plt.pause(0.3)                                # L466: yield to GUI event loop
\end{lstlisting}

\noindent \code{FuncAnimation} calls \code{update\_plot} every 300\,ms in the
matplotlib GUI event loop. The SLAM processing runs in the main thread;
visualisation runs in a \code{daemon} thread (dies when the main thread ends).

\section{Main Run Loop}

The main processing sequence (lines 269--348):

\begin{algorithm}[H]
\caption{IntegratedStereoSLAMSystem.run()}\label{alg:main}
\begin{algorithmic}[1]
\State Connect cameras; start viz thread
\While{running}
  \State $b \leftarrow$ \texttt{battery\_manager.update()}
  \If{\texttt{should\_skip\_frame(frame\_counter)}}
    \State \textbf{continue}                \Comment{save energy}
  \EndIf
  \State $(I_L, I_R) \leftarrow$ \texttt{camera\_manager.get\_stereo\_pair()}
  \State $d_{\min} \leftarrow 0.05 \times m_\text{batt}$   \Comment{adaptive threshold}
  \State \texttt{slam.min\_displacement} $\leftarrow d_{\min}$
  \State $\text{kf} \leftarrow$ \texttt{slam.process\_stereo\_pair}$(I_L, I_R)$
  \If{$N_\text{kf} > N_\text{last\_ba} + \Delta_\text{ba}$}
    \State \texttt{slam.optimize\_bundle()}
  \EndIf
  \State \texttt{display()} and handle keyboard
  \State compute FPS
\EndWhile
\State \texttt{cleanup(); save\_map()}
\end{algorithmic}
\end{algorithm}

%=============================================================================
\chapter{Module: \texttt{integrated\_slam\_rl.py}}
%=============================================================================

\section{Purpose}

This is the \emph{ultimate} entry point combining all sub-systems:
SLAM + RL + 2-D trajectory overlay + battery management + rich GUI.

\section{RL Integration Loop (Key Additions)}

\begin{lstlisting}
# Lines 299-330

# 1. Estimate state
uncertainty = rl_agent.uncertainty_metrics.update_from_slam(slam)  # L301
energy_need = energy_model.estimate_energy_reduction(battery_mode)  # L305

# 2. Get RL decision
rl_action = rl_agent.select_action(
    uncertainty, energy_need, battery_mode, training=training_mode)  # L311-313
should_create_keyframe = (rl_action == CREATE_KEYFRAME)              # L315

# 3. Learn in training mode
if training_mode:
    next_uncertainty = uncertainty + np.random.normal(0, 0.05)  # L320
    rl_agent.learn_from_experience(
        uncertainty, energy_need, battery_mode,
        prev_uncertainty, should_create_keyframe,
        next_uncertainty, energy_need)   # L323-327
\end{lstlisting}

\noindent In training mode the agent learns \emph{online} while the SLAM system
runs. The ``next uncertainty'' is approximated by adding Gaussian noise to the
current uncertainty (line 320), simulating the uncertainty dynamics without
waiting one full frame.

\section{2-D Trajectory Overlay}

\begin{lstlisting}
# Lines 400-453
traj_2d = trajectory[:, :2]         # project to XY plane
scale = (overlay_size - 2*margin) / max(max_x - min_x, max_y - min_y)
traj_scaled = (traj_2d - [min_x, min_y]) * scale + margin
# Draw coloured line (gradient from blue to red over time)
for i in range(len(traj_scaled) - 1):
    color = (int(255*(i/len(traj_scaled))), 100,
             int(100 + 155*(i/len(traj_scaled))))
    cv2.line(overlay, pt1, pt2, color, 2)
# Blend onto left frame
left_resized[0:H, 0:W] = cv2.addWeighted(
    left_resized[0:H, 0:W], 0.6, overlay, 0.4, 0)
\end{lstlisting}

The colour gradient maps time to hue: early waypoints are blue
$(0, 100, 100)$, later ones trend toward red $(255, 100, 255)$.
\code{cv2.addWeighted} computes: $I_{\text{blend}} = 0.6\,I_{\text{cam}} + 0.4\,I_{\text{overlay}}$.

%=============================================================================
\chapter{Module: \texttt{train\_rl\_agent.py}}
%=============================================================================

\section{SLAM Simulator}

The simulator provides a fast, SLAM-free environment for offline RL training.

\begin{lstlisting}
class SLAMSimulator:
    def reset(self):
        self.uncertainty = np.random.uniform(0.2, 0.8)  # L28
        self.energy_need = np.random.uniform(0.1, 0.9)  # L29
        self.battery_mode = np.random.choice(
            ['critical','low','moderate','good'])         # L30

    def step(self, action: KeyframeAction):
        if action == CREATE_KEYFRAME:
            self.uncertainty -= np.random.uniform(0.05, 0.15)  # L54
            self.energy_need += np.random.uniform(0.05, 0.10)  # L58
        else:
            self.uncertainty += np.random.uniform(0.02, 0.08)  # L62
            self.energy_need -= np.random.uniform(0.05, 0.15)  # L65
        # Stochastic noise
        self.uncertainty += np.random.normal(0, 0.02)   # L69
        # Battery mode change (10% probability)
        if np.random.random() < 0.1:
            self.battery_mode = np.random.choice(...)   # L73
\end{lstlisting}

\begin{notebox}[Why Simulate Before Deployment?]
Running the SLAM pipeline for 500 episodes (50\,000 steps) requires 500 camera
sessions. The simulator enables training in seconds on a CPU with a well-defined
transition model, providing a pre-trained policy that requires only minimal
fine-tuning on the real system.
\end{notebox}

\section{Training Loop}

\begin{lstlisting}
for episode in range(episodes):       # L161
    state = env.reset()
    for step in range(episode_length):  # L167
        action = agent.select_action(
            state['uncertainty'], state['energy_need'],
            state['battery_mode'], training=True)         # L169-173
        next_state, reward, done, info = env.step(action) # L177
        agent.learn_from_experience(...)                   # L180-189
        state = next_state
    # Epsilon decay per episode
    agent.epsilon *= epsilon_decay           # L198: 0.995 per episode
    agent.epsilon = max(agent.epsilon, 0.01) # L199: floor at 1%
\end{lstlisting}

\begin{mathbox}[Epsilon Decay Schedule]
Exponential decay per episode with factor $\delta = 0.995$:
\[
  \varepsilon_n = \max\!\left(\varepsilon_0 \cdot \delta^n,\; \varepsilon_{\min}\right)
\]
After 500 episodes: $\varepsilon_{500} = 1.0 \times 0.995^{500} \approx 0.082$,
just above the floor of 0.01.
\end{mathbox}

%=============================================================================
\chapter{Tools and Launcher Scripts}
%=============================================================================

\section{Dynamic IP Configuration (\texttt{tools/dynamic\_ip\_config.py})}

Provides a GUI or command-line tool to discover and configure camera IP
addresses on the local network. Key functionality:
\begin{itemize}
  \item Network scan for ESP32-CAM / Raspberry Pi camera devices
  \item Saves discovered IPs to a JSON configuration file
  \item Provides a Tkinter GUI for easy configuration
\end{itemize}

\section{Launcher Scripts Overview}

\begin{table}[H]
\centering
\caption{Launcher script purposes}\label{tab:launchers}
\begin{tabular}{lp{9cm}}
\toprule
\textbf{Script} & \textbf{Launches} \\
\midrule
\file{launch\_basic\_slam.py} & Basic stereo SLAM without RL \\
\file{launch\_integrated\_slam.py} & Battery-aware SLAM (\file{integrated\_slam.py}) \\
\file{launch\_integrated\_slam\_rl.py} & Full SLAM+RL system \\
\file{launch\_rl\_trainer.py} & Offline RL training on SLAM simulator \\
\file{launch\_rl\_gui.py} & Training dashboard GUI \\
\file{launch\_stereo\_slam.py} & Stereo-only SLAM \\
\file{play\_dataset\_video.py} & Replay saved SLAM map as video \\
\file{run\_backend.py} & Backend API server \\
\bottomrule
\end{tabular}
\end{table}

%=============================================================================
\chapter{Data Flow: End-to-End Trace}
%=============================================================================

\section{Single Frame Processing Trace}

\begin{enumerate}[leftmargin=1.5em, label=\textbf{Step~\arabic*.}]

  \item \textbf{Frame acquisition} (\file{camera\_stream\_handler.py})\\
        HTTP chunks arrive $\to$ JPEG boundaries detected $\to$ \code{cv2.imdecode}
        $\to$ NumPy BGR array \code{(H,W,3)}.

  \item \textbf{Battery check} (\file{integrated\_slam.py} or \file{integrated\_slam\_rl.py})\\
        \code{psutil.sensors\_battery()} queried every 5\,s $\to$ mode
        \{critical, low, moderate, good\} $\to$ frame-skip decision.

  \item \textbf{Feature extraction} (\file{advanced\_features.py})\\
        CLAHE $\to$ ORB/SIFT/AKAZE \code{detectAndCompute}
        $\to$ \code{(keypoints, descriptors, quality\_score)}.

  \item \textbf{Stereo matching} (\file{advanced\_features.py})\\
        BF-matcher \code{knnMatch(k=2)} $\to$ Lowe ratio test $\to$
        list of \code{(left\_idx, right\_idx)} pairs.

  \item \textbf{Depth estimation} (\file{stereo\_vision.py})\\
        \code{StereoBM.compute} $\to$ disparity map $\to$
        $Z = fb/d$ per matched pixel.

  \item \textbf{3-D triangulation} (\file{stereo\_vision.py})\\
        Back-projection: $(X,Y,Z) = (u-c_x,\, v-c_y,\, f) \cdot Z/f$.

  \item \textbf{Visual odometry} (\file{visual\_odometry.py})\\
        \code{findEssentialMat} (RANSAC) $\to$ \code{recoverPose}
        $\to$ $(\mat{R}, \vect{t})$ $\to$ accumulated pose $\mat{T}$.

  \item \textbf{RL decision} (\file{rl\_keyframe\_selector.py} or \file{dqn\_agent.py})\\
        Compute uncertainty $u$ and energy need $e$ $\to$
        $\varepsilon$-greedy $Q$-table / neural network lookup
        $\to$ \{SKIP, CREATE\}.

  \item \textbf{Keyframe creation} (\file{stereo\_slam.py})\\
        If CREATE: store pose + descriptors + 3-D points $\to$ loop closure
        check $\to$ (periodic) bundle adjustment.

  \item \textbf{Q-value update} (training mode)\\
        Next-state uncertainty computed $\to$ TD error $\delta$
        $\to$ $Q \leftarrow Q + \alpha\delta$.

  \item \textbf{Visualisation}\\
        Point cloud + trajectory rendered in 3-D Matplotlib window.
        2-D trajectory overlay blended onto left camera feed.
        Statistics displayed via OpenCV \code{putText}.

  \item \textbf{Persistence}\\
        On quit: \code{save\_map()} serialises point cloud + trajectory to JSON.
        (Optionally: PLY export for MeshLab / CloudCompare.)
\end{enumerate}

%=============================================================================
\chapter{Configuration and Hyper-Parameters}
%=============================================================================

\begin{longtable}{|p{4.5cm}|p{2.5cm}|p{6.5cm}|}
\caption{Complete hyper-parameter reference}\label{tab:hyperparams}\\
\hline
\textbf{Parameter} & \textbf{Default} & \textbf{Description / Effect} \\ \hline
\endfirsthead
\hline
\textbf{Parameter} & \textbf{Default} & \textbf{Description / Effect} \\ \hline
\endhead
\multicolumn{3}{c}{\textit{Camera \& Stereo}} \\ \hline
\code{focal\_length} & 500 px & Pinhole camera focal length \\ \hline
\code{baseline} & 0.12 m & Stereo camera separation \\ \hline
\code{numDisparities} & 128 & StereoBM search range \\ \hline
\code{blockSize} & 15 & StereoBM correlation window \\ \hline
\multicolumn{3}{c}{\textit{SLAM}} \\ \hline
\code{max\_keyframes} & 100 & Max keyframes in circular buffer \\ \hline
\code{min\_displacement} & 0.05 m & Min translation for new keyframe \\ \hline
\code{loop\_closure\_threshold} & 0.75 & Min visual match score for loop closure \\ \hline
\code{max\_features} (ORB) & 500 & Max ORB keypoints per frame \\ \hline
\multicolumn{3}{c}{\textit{Q-Learning Agent}} \\ \hline
\code{learning\_rate} $\alpha$ & 0.1 & Q-value update step size \\ \hline
\code{discount\_factor} $\gamma$ & 0.95 & Future reward discount \\ \hline
\code{epsilon} & 0.1 & Exploration rate \\ \hline
\code{state\_bins} & 10 & Discretisation resolution \\ \hline
\multicolumn{3}{c}{\textit{DQN Agent}} \\ \hline
\code{hidden\_dim} & 128 & Neural network hidden size \\ \hline
\code{learning\_rate} & 1e-4 & Adam optimiser LR \\ \hline
\code{gamma} & 0.99 & Discount factor \\ \hline
\code{epsilon\_start} & 0.3 & Initial exploration \\ \hline
\code{epsilon\_end} & 0.05 & Final exploration \\ \hline
\code{epsilon\_decay\_steps} & 100\,000 & Steps to decay over \\ \hline
\code{buffer capacity} & 100\,000 & Replay buffer size \\ \hline
\code{batch\_size} & 64 & Training batch size \\ \hline
\code{target\_update\_freq} & 1000 & Steps between target net sync \\ \hline
\multicolumn{3}{c}{\textit{Reward Weights}} \\ \hline
\code{W\_BENEFIT} & 8.0 & Localisation improvement weight \\ \hline
\code{W\_ENERGY} & 6.0 & Energy cost weight \\ \hline
\code{W\_TRACKING\_LOSS} & 40.0 & Tracking failure penalty \\ \hline
\code{BONUS\_SKIP} & 2.0 & Successful skip bonus \\ \hline
\code{PENALTY\_WASTE} & 3.0 & Wasted keyframe penalty \\ \hline
\multicolumn{3}{c}{\textit{Battery Thresholds}} \\ \hline
\code{BATTERY\_CRITICAL} & 5\% & Process 1/8 frames \\ \hline
\code{BATTERY\_LOW} & 15\% & Process 1/4 frames \\ \hline
\code{BATTERY\_NORMAL} & 50\% & Process 1/2 frames \\ \hline
\end{longtable}

%=============================================================================
\chapter{Algorithms Summary}
%=============================================================================

\begin{table}[H]
\centering
\caption{Key algorithms and their complexity}\label{tab:algos}
\begin{tabular}{llll}
\toprule
\textbf{Algorithm} & \textbf{Location} & \textbf{Time} & \textbf{Notes} \\
\midrule
MJPEG boundary scan & \file{camera\_stream\_handler} & $\mathcal{O}(C)$ & $C$ = chunk size \\
CLAHE & \file{advanced\_features} & $\mathcal{O}(HW)$ & per frame \\
SIFT detection & \file{stereo\_vision} & $\mathcal{O}(HW)$ & scale space \\
ORB detection & \file{advanced\_features} & $\mathcal{O}(HW)$ & faster than SIFT \\
BF kNN matching & \file{advanced\_features} & $\mathcal{O}(N_1 N_2 D)$ & $D$=descriptor dim \\
Lowe ratio test & \file{stereo\_vision} & $\mathcal{O}(M)$ & $M$ = match count \\
StereoBM & \file{stereo\_vision} & $\mathcal{O}(HWd_{\max})$ & $d_{\max}$=128 \\
Back-projection & \file{stereo\_vision} & $\mathcal{O}(M)$ & vectorised \\
Essential matrix RANSAC & \file{visual\_odometry} & $\mathcal{O}(N\log N)$ & \\
Kabsch SVD & \file{stereo\_vision} & $\mathcal{O}(N)$ & SVD of $3\times3$ \\
Loop closure check & \file{stereo\_slam} & $\mathcal{O}(K \cdot N_1 N_2 D)$ & $K$ = old keyframes \\
Voxel downsampling & \file{advanced\_features} & $\mathcal{O}(N\log N)$ & \\
Q-table lookup & \file{rl\_keyframe\_selector} & $\mathcal{O}(1)$ & hash map \\
DQN forward pass & \file{dqn\_agent} & $\mathcal{O}(7{\cdot}128 + 128^2 + 128{\cdot}3)$ & tiny net \\
Double DQN update & \file{dqn\_agent} & $\mathcal{O}(B \cdot \text{net})$ & $B=64$ \\
\bottomrule
\end{tabular}
\end{table}

%=============================================================================
\chapter{Known Limitations and Future Work}
%=============================================================================

\section{Current Limitations}

\begin{enumerate}
  \item \textbf{No hardware synchronisation}: The two cameras are polled
        independently; temporal mismatch can introduce small but systematic
        errors in stereo depth estimation.

  \item \textbf{Fixed intrinsics}: Focal length (500 px) and principal point
        ($320, 240$) are hardcoded estimates. A proper calibration step using a
        checkerboard and \code{cv2.calibrateCamera} would substantially improve
        3-D accuracy.

  \item \textbf{Simplified bundle adjustment}: The \code{optimize\_bundle}
        method is a placeholder. Full BA requires a non-linear least-squares
        solver (e.g., Ceres, g2o) minimising reprojection error over all
        poses and points simultaneously.

  \item \textbf{No pose-graph optimisation}: Loop closures are detected but their
        constraints are not used to close the trajectory loop mathematically.

  \item \textbf{Tabular Q-table scalability}: The Q-table has $400$ states;
        for richer state representations (e.g., including optical flow, IMU) the
        state space explodes. The DQN agent addresses this but is not yet
        integrated into the main pipeline.

  \item \textbf{Online RL noise}: The next-state uncertainty approximation
        (Gaussian noise) in training mode is inaccurate; a proper transition
        model or off-policy training would be more principled.
\end{enumerate}

\section{Future Extensions}

\begin{enumerate}
  \item \textbf{IMU fusion}: Integrate accelerometer/gyroscope data to improve
        pose estimation between frames, especially under fast motion or poor
        lighting.

  \item \textbf{Semantic SLAM}: Use a small object-detection network (e.g., YOLO
        nano) to label landmarks semantically, enabling loop closure via
        object co-occurrence.

  \item \textbf{Cloud RL Training}: Offload DQN training to a GPU server in
        the cloud; the robot downloads the latest policy weights periodically
        (federated RL).

  \item \textbf{Pose-Graph SLAM}: Integrate iSAM2 or GTSAM for proper maximum a
        posteriori pose estimation over the full trajectory.

  \item \textbf{Curriculum Training}: Start the SLAM simulator with easy episodes
        (high battery, low uncertainty) and progressively increase difficulty,
        a technique shown to accelerate RL convergence substantially.
\end{enumerate}

%=============================================================================
\chapter{Glossary}
%=============================================================================

\begin{description}[leftmargin=3cm, style=nextline]

  \item[BA] Bundle Adjustment --- joint optimisation of camera poses and 3-D
       point positions.

  \item[CLAHE] Contrast-Limited Adaptive Histogram Equalization.

  \item[DQN] Deep Q-Network --- Q-learning with a neural network function
       approximator.

  \item[DoG] Difference of Gaussians --- scale-space derivative used in SIFT.

  \item[ε-greedy] Exploration policy: random with probability ε, greedy otherwise.

  \item[Essential Matrix] $3\times3$ matrix encoding relative rotation and
       translation between two calibrated camera views.

  \item[Fundamental Matrix] $3\times3$ matrix encoding epipolar geometry for
       uncalibrated cameras.

  \item[IMU] Inertial Measurement Unit --- accelerometer + gyroscope.

  \item[Kabsch] SVD-based algorithm for optimal rigid-body alignment.

  \item[Keyframe] A selected image frame whose pose and features are stored in
       the SLAM map.

  \item[Loop Closure] Detection that the camera has returned to a previously
       visited location, used to correct accumulated drift.

  \item[MDP] Markov Decision Process.

  \item[MJPEG] Motion JPEG --- video stream as consecutive JPEG frames.

  \item[ORB] Oriented FAST and Rotated BRIEF --- fast binary feature detector.

  \item[Pose] A $4\times4$ SE(3) rigid-body transformation matrix.

  \item[RANSAC] Random Sample Consensus --- robust estimation in the presence
       of outliers.

  \item[Reward] Scalar feedback signal in RL indicating action quality.

  \item[SIFT] Scale-Invariant Feature Transform (Lowe, 2004).

  \item[SLAM] Simultaneous Localisation and Mapping.

  \item[SE(3)] Special Euclidean group of 3-D rigid transformations.

  \item[SO(3)] Special Orthogonal group of 3-D rotations.

  \item[State Space] Discretised $10\times10\times4 = 400$ states encoding
       uncertainty, energy need, and battery mode.

  \item[StereoBM] OpenCV's block-matching stereo disparity algorithm.

  \item[TD Error] Temporal-Difference error: $\delta = r + \gamma\max Q(s',a') - Q(s,a)$.

  \item[Trajectory] Sequence of camera centre positions in world coordinates.

  \item[Triangulation] Computing 3-D point position from 2-D stereo matches.

  \item[Voxel Downsampling] Reducing point cloud density via spatial binning.
\end{description}

%=============================================================================
\appendix
%=============================================================================

\chapter{Complete Class Hierarchy}

\begin{longtable}{|p{5cm}|p{3cm}|p{5.5cm}|}
\caption{All classes and their module}\label{tab:classes}\\
\hline
\textbf{Class} & \textbf{Module} & \textbf{Key Methods} \\ \hline
\endfirsthead
\hline
\textbf{Class} & \textbf{Module} & \textbf{Key Methods} \\ \hline
\endhead
\code{CameraStreamHandler} & camera\_stream\_handler & start, stop, get\_frame \\ \hline
\code{StereoCameraManager} & camera\_stream\_handler & start, stop, get\_stereo\_pair \\ \hline
\code{StereoCalibratedPipeline} & stereo\_vision & extract\_features, match\_features, estimate\_depth, compute\_3d\_points \\ \hline
\code{AdvancedFeatureExtractor} & advanced\_features & extract\_features \\ \hline
\code{RobustStereoMatcher} & advanced\_features & match\_stereo, match\_temporal \\ \hline
\code{EpipolarGeometry} & advanced\_features & compute\_fundamental\_matrix, filter\_by\_epipolar\_constraint \\ \hline
\code{PointCloudFilter} & advanced\_features & remove\_outliers, voxel\_grid\_downsample, remove\_duplicates \\ \hline
\code{LoopClosureDetector} & advanced\_features & add\_pose, detect\_loop\_closure \\ \hline
\code{BundleAdjustment} & advanced\_features & optimize\_poses \\ \hline
\code{CloudExporter} & advanced\_features & export\_ply, export\_trajectory\_json \\ \hline
\code{VisualOdometryEngine} & visual\_odometry & process\_stereo\_pair, get\_trajectory, get\_camera\_pose \\ \hline
\code{KeyFrame} & stereo\_slam & get\_camera\_center \\ \hline
\code{StereoSLAMSystem} & stereo\_slam & process\_stereo\_pair, optimize\_bundle, save\_map \\ \hline
\code{KeyframeAction} & rl\_keyframe\_selector & SKIP\_FRAME, CREATE\_KEYFRAME \\ \hline
\code{LocalizationUncertaintyMetrics} & rl\_keyframe\_selector & update\_from\_slam \\ \hline
\code{EnergyModel} & rl\_keyframe\_selector & estimate\_energy\_reduction \\ \hline
\code{RLKeyframeSelector} & rl\_keyframe\_selector & select\_action, compute\_reward, update\_q\_value, learn\_from\_experience \\ \hline
\code{RL\_SLAM\_Integration} & rl\_keyframe\_selector & decide\_keyframe\_creation \\ \hline
\code{DQNNetwork} & dqn\_agent & forward \\ \hline
\code{ReplayBuffer} & dqn\_agent & push, sample \\ \hline
\code{DQNAgent} & dqn\_agent & select\_action, train\_step, save\_model, load\_model \\ \hline
\code{EnergyModel} & energy\_battery\_model & get\_action\_cost, estimate\_remaining\_runtime \\ \hline
\code{BatteryMonitor} & energy\_battery\_model & update, get\_degradation\_level, get\_energy\_weight\_multiplier \\ \hline
\code{EnergyAwareRewardCalculator} & energy\_battery\_model & compute\_reward, compute\_batch\_reward \\ \hline
\code{DegradationLevel} & degradation\_policy & (data class) \\ \hline
\code{SLAMDegradationPolicy} & degradation\_policy & update\_battery, should\_create\_keyframe, subsample\_descriptors \\ \hline
\code{SLAMSimulator} & train\_rl\_agent & reset, step \\ \hline
\code{BatteryManager} & integrated\_slam & update\_battery\_state, get\_mode, should\_skip\_frame \\ \hline
\code{IntegratedSLAMVisualizer} & integrated\_slam & initialize, render \\ \hline
\code{IntegratedStereoSLAMSystem} & integrated\_slam & run, cleanup \\ \hline
\code{UltimateIntegratedStereoSLAMWithRL} & integrated\_slam\_rl & run, cleanup \\ \hline
\end{longtable}

\chapter{Quick-Start Commands}

\begin{lstlisting}[language=bash, style=pythonstyle]
# Install dependencies
pip install opencv-python numpy torch matplotlib psutil requests scipy

# Run basic SLAM (no RL)
python integrated_slam.py --left-camera http://192.168.1.47/stream \
                           --right-camera http://192.168.1.48/stream

# Run SLAM + RL (inference mode, pre-trained policy)
python app/backend/integrated_slam_rl.py \
    --left-camera http://192.168.1.47/stream \
    --right-camera http://192.168.1.48/stream \
    --model rl_keyframe_model.json

# Train RL agent offline (500 episodes)
python app/backend/train_rl_agent.py \
    --episodes 500 --length 100 --lr 0.1 --output rl_keyframe_model.json

# Run SLAM + RL (online training mode)
python app/backend/integrated_slam_rl.py \
    --left-camera http://192.168.1.47/stream \
    --right-camera http://192.168.1.48/stream \
    --train --model rl_keyframe_model.json
\end{lstlisting}

\chapter{Mathematical Reference Sheet}

\begin{mathbox}[Key Equations Summary]
\textbf{Stereo Depth:} $Z = \dfrac{f \cdot b}{d}$ \\[6pt]
\textbf{Back-projection:} $X = \dfrac{(u-c_x)Z}{f},\; Y = \dfrac{(v-c_y)Z}{f}$ \\[6pt]
\textbf{Epipolar:} $\vect{x}_2^\T \mat{E}\, \vect{x}_1 = 0$ \\[6pt]
\textbf{Pose update:} $\mat{T}_{k+1} = \mat{T}_k \cdot \mat{T}_\Delta$ \\[6pt]
\textbf{Q-learning:} $Q(s,a) \leftarrow Q(s,a) + \alpha\bigl[r + \gamma\max_{a'}Q(s',a') - Q(s,a)\bigr]$ \\[6pt]
\textbf{Double DQN target:} $y = r + \gamma Q_{\theta^-}\!\bigl(s',\argmax_{a'}Q_\theta(s',a')\bigr)$ \\[6pt]
\textbf{ε-greedy:} $\pi(s) = \begin{cases}\text{random} & \text{prob } \varepsilon \\ \argmax_a Q(s,a) & \text{prob } 1-\varepsilon\end{cases}$ \\[6pt]
\textbf{Energy weight:} $w_e = 1 + k_b(1-b/100)$ \\[6pt]
\textbf{Spatial entropy:} $H = -\sum_k p_k \ln p_k$ \\[6pt]
\textbf{Kabsch:} $\mat{R}^* = \mat{V}\,\mathrm{diag}(1,1,\det(\mat{VU}^\T))\,\mat{U}^\T$ \\[6pt]
\textbf{Lowe ratio:} $\norm{\vect{d}-\vect{d}^1} / \norm{\vect{d}-\vect{d}^2} < 0.7$ \\[6pt]
\textbf{Gradient clipping:} $\nabla \leftarrow \nabla / \max(1, \norm{\nabla}/g_{\max})$
\end{mathbox}

%----------------------------- BIBLIOGRAPHY --------------------------------
\begin{thebibliography}{99}

\bibitem{lowe2004}
D. G. Lowe,
``Distinctive image features from scale-invariant keypoints,''
\textit{International Journal of Computer Vision}, vol.~60, no.~2, pp.~91--110, 2004.

\bibitem{mnih2015}
V. Mnih et al.,
``Human-level control through deep reinforcement learning,''
\textit{Nature}, vol.~518, pp.~529--533, 2015.

\bibitem{hasselt2016}
H. van Hasselt, A. Guez, and D. Silver,
``Deep reinforcement learning with double Q-learning,''
\textit{Proc. AAAI}, 2016.

\bibitem{kabsch1976}
W. Kabsch,
``A solution for the best rotation to relate two sets of vectors,''
\textit{Acta Crystallographica A}, vol.~32, pp.~922--923, 1976.

\bibitem{hartley2004}
R. Hartley and A. Zisserman,
\textit{Multiple View Geometry in Computer Vision}, 2nd~ed.
Cambridge University Press, 2004.

\bibitem{thrun2005}
S. Thrun, W. Burgard, and D. Fox,
\textit{Probabilistic Robotics}.
MIT Press, 2005.

\bibitem{sutton2018}
R. S. Sutton and A. G. Barto,
\textit{Reinforcement Learning: An Introduction}, 2nd~ed.
MIT Press, 2018.

\bibitem{rublee2011}
E. Rublee et al.,
``ORB: An efficient alternative to SIFT or SURF,''
\textit{Proc. ICCV}, 2011.

\bibitem{watkins1992}
C. J. C. H. Watkins and P. Dayan,
``Q-learning,''
\textit{Machine Learning}, vol.~8, pp.~279--292, 1992.

\end{thebibliography}

\end{document}
